\documentclass[11pt]{article}
\usepackage[margin=1in]{geometry}
\usepackage{amsmath,amssymb,amsthm}
\usepackage[T1]{fontenc}
\usepackage{lmodern}
\usepackage[hidelinks]{hyperref}
\newtheorem{theorem}{Theorem}
\newtheorem{lemma}{Lemma}
\newcommand{\Z}{\mathbb Z}
\newcommand{\N}{\mathbb N}
\title{Prime Radicands Contradict the Printed Density\newline
  in Conjecture 00000000259}
\author{idealistichacker\\\small AI-assisted submission prepared for review}
\date{October 22, 2026}

\begin{document}
\maketitle
\begin{abstract}
The printed conjecture concerns primes $p\equiv1\pmod4$, not general composite
squarefree radicands. Every such prime admits an integer solution of
$x^2-py^2=-1$. Consequently the success proportion within that prime class is
identically one once the counting denominator is nonzero. The printed constant
$1-\prod_{k\ge2}(1-2^{-k})$ belongs to $[1/4,7/16]$ and is not one. We give the
descent proof and an actual finite-count/limit formalization. This draft does
not claim to resolve a different negative-Pell density problem for composite
discriminants, and is not a claim of publication or upstream acceptance.
\end{abstract}

\section{Exact Statement and Interpretation}
The English source asserts: ``The density of primes $p\equiv1\pmod4$ for which
the negative Pell equation is solvable equals
$1-\prod_{k\ge2}(1-2^{-k})$.'' The Chinese source likewise explicitly specifies
primes. Both complete source languages are locked at:
\begin{quote}\small\ttfamily
Git blob: 15fe4da518462416cfbace31ff72247830e4ea6b\\
Commit: b355e42906e6eaf52bc7e9d005e0ef02d8aa7400.
\end{quote}

Here density within the displayed reference class means the standard relative
counting density. Define
\[
 B_N=\{p<N:p\text{ prime},\ p\equiv1\pmod4\},\qquad
 A_N=\{p\in B_N:\exists x,y\in\Z,\ x^2-py^2=-1\}.
\]
The assertion formalized and refuted is
\[
 \frac{|A_N|}{|B_N|}\longrightarrow
 C:=1-\prod_{k\ge2}(1-2^{-k})\quad(N\longrightarrow\infty).
\]
The denominator is explicitly disclosed; the Lean theorem is not advertised as
a proof about a different density denominator. No restriction to primes with
an additional class-number condition is inserted. A prime $p\equiv1\pmod4$
is already a squarefree positive fundamental discriminant. The introductory
wide-class-number assertion is not needed as an assumption: one false necessary
density clause suffices to refute the displayed conjunction under this reading.

\section{All Eligible Primes Give Negative Pell Solutions}
\begin{theorem}
For every prime $p\equiv1\pmod4$, there exist $u,v\in\Z$ with
$u^2-pv^2=-1$.
\end{theorem}
\begin{proof}
The ordinary Pell existence theorem supplies an integer solution
$x^2-py^2=1$ whose $x>1$ is minimal among positive Pell solutions with
$x>1$. This is applicable because a positive prime is not a square.
Reducing modulo four shows $x$ odd and $y$ even. Set
\[
 a=(x-1)/2>0,\qquad b=(x+1)/2=a+1,\qquad c=y/2.
\]
Then $a,b$ are coprime positive integers and $ab=pc^2$. Unique factorization
and coprimality put the prime $p$ in exactly one factor. Each remaining prime
exponent is even, so either
\[
 a=pu^2,\quad b=v^2,\qquad\text{or}\qquad
 a=u^2,\quad b=pv^2.
\]
In the first case $v^2-pu^2=b-a=1$. Since $b=a+1>1$ and
$v^2=(x+1)/2$, we have $1<|v|<x$. Thus $(|v|,|u|)$ is a positive Pell
solution with smaller first coordinate, contradicting minimality. The second
case therefore holds, giving $u^2-pv^2=a-b=-1$.
\end{proof}

The Lean proof invokes Mathlib's actual
\texttt{Pell.IsFundamental.exists\_of\_not\_isSquare}; it does not assume
the required negative-Pell solution. Its descent uses actual integer Pell
solutions and a proved split of adjacent coprime square factors. As a diagnostic
example, $p=109$ has the exact certificate
\[
 8890182^2-109\cdot851525^2=-1.
\]
A bounded search that stops before $y=851525$ is not a proof of insolubility.
The general argument above, not this isolated certificate, proves the theorem.

\section{The Printed Constant Is Separated From One}
Write $f_i=1-2^{-(i+2)}$ for $i\in\N$ and $P_n=\prod_{i<n}f_i$.
The geometric series of deviations is summable, so $\prod_i f_i$ exists.
All factors lie in $[0,1]$.

For $0\le a_i\le1$, the elementary union bound for products states
\[
 \prod_{i<n}(1-a_i)\ge1-\sum_{i<n}a_i.
\]
It follows by induction: if $Q\le1$ and $Q\ge1-S$, then
$Q(1-a)=Q-aQ\ge Q-a\ge1-S-a$. The shifted geometric tail satisfies
$\sum_{i<n}2^{-(i+3)}\le1/4$. Splitting off the first factor gives
\[
 \frac9{16}\le P_{n+1}
 =\frac34\prod_{i<n}(1-2^{-(i+3)})\le\frac34.
\]
Passing to the genuine infinite-product limit yields
\[
 \frac9{16}\le\prod_{k\ge2}(1-2^{-k})\le\frac34,
 \qquad \frac14\le C\le\frac7{16}<1.
\]
These are exact rational inequalities, not numerical approximations to an
uncomputed product. Lean proves summability, multipliability, finite-product
bounds, and the passage to the limit.

\section{Counting and Closed Limit Negation}
\begin{theorem}
$A_N=B_N$ for every $N\in\N$. For every $N\ge6$,
$|A_N|/|B_N|=1$. In particular this ratio cannot tend to $C$.
\end{theorem}
\begin{proof}
The prime-radicand theorem applies to each member of $B_N$, so $A_N=B_N$.
For $N\ge6$, $5\in B_N$, making $|B_N|>0$. Dividing equal counts gives
one. The ratio is therefore eventually the constant one and converges to one.
Uniqueness of real limits and $C\ne1$ negate convergence to the conjectured
value. No prime number theorem or statistical heuristic is required for this
relative-count identity.
\end{proof}

\section{Formalization and Verification Boundary}
The research sources pin Lean 4.33.1. The Mathlib commit is:
\begin{quote}\small\ttfamily
0df444a360eaa60ab8c11dca51a86af692955474.
\end{quote}
\begin{itemize}
\item \texttt{PrimePell.lean}: four general integer Pell theorems.
\item \texttt{ProductBounds.lean}: eight actual infinite-product and bound theorems.
\item \texttt{RelativeDensity.lean}: seven count, limit, negation, and certificate
  theorems, with the actual sets $A_N,B_N$ defined using finite-set filters.
\item \texttt{Main.lean}: imports the complete mathematical argument.
\item \texttt{Check.lean}: all 19 actual theorem types and axiom audits.
\end{itemize}
The decisive declaration is
\begin{quote}\small\ttfamily
PellDensity259.not\_claimed\_relative\_density :\newline
\phantom{xx}not (Filter.Tendsto relativeDensity Filter.atTop\newline
\phantom{xxxx}(nhds PellConstant259.claimedConstant)).
\end{quote}
It has no external density, solution-existence, or product-separation premise.
Local warning-as-error source builds and all axiom audits have succeeded; only
the standard foundations \texttt{propext}, \texttt{Classical.choice}, and
\texttt{Quot.sound} occur. The package reproducer verifies all nine physical
pinned Git dependencies, their origins, tracked source bytes, and index/tree
agreement before and after execution. It cleans only this project's build
outputs, rebuilds Main, replays every mathematical source and Main with warnings
as errors, and audits all 19 actual types and axiom footprints. Its exact integer
certificates and finite product checks are regression diagnostics, not replacements
for the Lean proofs. Installed Lean and dependency build caches remain trust
inputs: this is not an empty-machine bootstrap or a clean build of all Mathlib
dependencies. Independent review is internal AI-assisted review, not human,
institutional, or organizer endorsement. Submission requests upstream review;
it is not a claim of acceptance or first-solver credit.
\end{document}
